\chapter{From integral modules to complexes and stable invariants}
\label{ch:orders-complexes}

\section{What the second volume contributes}

The second supplied volume is Klaus W. Roggenkamp's \emph{Lattices over Orders II}, published in 1970 as volume 142 of the Lecture Notes in Mathematics. Its organizing questions concern integral modules over orders: how local isomorphisms differ from global isomorphisms, how genera are classified, when cancellation is valid, and when only finitely many indecomposable lattices occur. Chapters VI and VII develop modules, ideals and genera; Chapter VIII constructs Grothendieck and Whitehead groups; Chapters IX and X study hereditary and Bass orders and questions of finite lattice type. These subjects supply a precise language for passing from individual arithmetic objects to stable invariants. They do not identify cryptographic hardness with a topological index or an analytic spectral gap.\cite{v2Roggenkamp1970II}

There is a particularly direct historical connection to the analytic side of this monograph. The remark following VIII.2.4, printed page 114, discusses Atiyah--Hirzebruch topological K-theory and the correspondence between vector bundles and projective modules. This is an actual connection in the supplied author's text, rather than an analogy inferred from a shared word. Its modern use requires distinguishing several groups that have similar notation. In this chapter these distinctions are part of the mathematical statement, not qualifications added after a conclusion.\cite[Chapter VIII, \S\S1--2]{v2Roggenkamp1970II}

All results below that have a displayed proof are derived within the stated framework. Deep classification and index results are explicitly identified as external theorems. The purpose is to expose the exact points at which an arithmetic construction can be transported to a complex, a complex can be transported to a stable class, and a stable class can be evaluated by an index. Each transport needs data and hypotheses. A chain of related subjects is not itself a theorem connecting their computational difficulties.

\section{The exact category of lattices}

Let $R$ be a Dedekind domain, let $K$ be its fraction field, and let $A$ be a finite-dimensional separable $K$-algebra. An $R$-order $\Lambda\subseteq A$ is a unital subring, finitely generated as an $R$-module, such that $K\Lambda=A$. A left $\Lambda$-lattice is a left $\Lambda$-module $M$ that is finitely generated and torsion-free over $R$. Because $R$ is Dedekind, this is equivalent to $M$ being finitely generated and projective over $R$. Projectivity over $R$ must not be confused with projectivity over $\Lambda$. The distinction is responsible for much of the homological content of the theory.\cite{v2Roggenkamp1970II}

\begin{proposition}[Extension closure of lattices]
Suppose
\[
0\longrightarrow M'\xrightarrow{u}M\xrightarrow{v}M''\longrightarrow0
\]
is an exact sequence of finitely generated $\Lambda$-modules. If $M'$ and $M''$ are $\Lambda$-lattices, then $M$ is a $\Lambda$-lattice. The sequence splits as a sequence of $R$-modules, although it need not split as a sequence of $\Lambda$-modules.
\end{proposition}
\begin{proof}
Finite generation over $R$ follows by taking generators of $M'$, lifting generators of $M''$, and using exactness. If $r\in R\setminus\{0\}$ and $rm=0$, then $r v(m)=0$ in the torsion-free module $M''$, so $v(m)=0$. Write $m=u(m')$. Injectivity of $u$ gives $rm'=0$, hence $m'=0$ and $m=0$. Thus $M$ is torsion-free. Finally, projectivity of $M''$ over $R$ supplies an $R$-linear section of $v$. Nothing in this argument produces a $\Lambda$-linear section.
\end{proof}

The lattice category therefore has a natural exact structure: its admissible exact sequences are those short sequences that are exact as module sequences and whose three terms are lattices. This permits extension classes without requiring the category itself to contain every kernel and cokernel of every morphism. For instance, multiplication by a nonunit on a nonzero lattice can have a torsion cokernel, which is not a lattice.

\begin{example}[Why quotient conditions matter]
The map $\mathbb Z\xrightarrow{\,p\,}\mathbb Z$, for a prime $p$, has cokernel $\mathbb Z/p\mathbb Z$. Both its source and target are $\mathbb Z$-lattices, while its cokernel is not. It is therefore incorrect to regard all maps between lattices as producing exact sequences inside the lattice category. The ambient module category and the exact lattice category encode different information.
\end{example}

An extension of $M''$ by $M'$ is an exact sequence as above, considered up to an isomorphism that is the identity on the two end terms. Its equivalence class lies in $\operatorname{Ext}^1_\Lambda(M'',M')$. When both endpoints are lattices, the preceding proposition shows that every such module extension has a lattice as middle term. Thus ordinary module $\operatorname{Ext}^1$ really does govern extensions within this exact category. The zero class is precisely the class of split extensions. This gives an initial, concrete bridge between the arithmetic category in Roggenkamp's books and the deformation calculations developed later.

\section{Three different Grothendieck groups}

For an exact category $\mathcal E$, write $G_0(\mathcal E)$ for the abelian group generated by isomorphism classes $[M]$ of objects, subject to
\[
[M]=[M']+[M'']
\]
whenever $0\to M'\to M\to M''\to0$ is admissibly exact. For the same additive category one can instead impose only direct-sum relations. Denote the resulting group by $K_0^{\mathrm{split}}(\mathcal E)$. There is a surjection
\[
K_0^{\mathrm{split}}(\mathcal E)\longrightarrow G_0(\mathcal E),
\]
because split exact sequences are exact sequences. In general it is not an isomorphism. Roggenkamp makes this direct-sum versus short-exact-sequence distinction at the beginning of Chapter VIII.\cite[Definitions VIII.1.2 and VIII.1.6]{v2Roggenkamp1970II}

For a unital ring $\Lambda$, the usual algebraic $K_0(\Lambda)$ is formed from finitely generated projective $\Lambda$-modules under direct sums. Since every short exact sequence of projective modules splits when its quotient is projective, the split and exact constructions agree on this category. In contrast,
\[
G_0^{\mathrm{lat}}(\Lambda)
  :=G_0(\operatorname{Lat}(\Lambda))
\]
uses all lattices and their exact sequences. The natural map $K_0(\Lambda)\to G_0^{\mathrm{lat}}(\Lambda)$ must not be declared an isomorphism without additional hypotheses.

\begin{proposition}[Universal property]
Let $B$ be an abelian group and let $a$ assign an element $a(M)\in B$ to every object of $\mathcal E$. Suppose $a$ is invariant under isomorphism and additive on admissible short exact sequences. There is a unique group homomorphism
\[
\overline a:G_0(\mathcal E)\longrightarrow B,
\qquad \overline a([M])=a(M).
\]
\end{proposition}
\begin{proof}
Map the free abelian group on isomorphism classes to $B$ by extending $a$ linearly. Every defining relation $[M]-[M']-[M'']$ maps to zero by additivity, so the map descends to the quotient. Its values on the generators prescribe it uniquely.
\end{proof}

This simple property is the proper explanation of the usefulness of Grothendieck groups. They collect additive invariants. They do not retain every invariant of an object, and they are not automatic certificates of isomorphism.

\begin{proposition}[Equality in a split group]
For objects $M,N$ in an additive category, equality of their classes in $K_0^{\mathrm{split}}$ holds if and only if there exists an object $T$ such that $M\oplus T\cong N\oplus T$.
\end{proposition}
\begin{proof}
The implication from a stable isomorphism follows directly from the relation for direct sums. For the reverse implication, form the commutative monoid of isomorphism classes under direct sum. Its group completion consists of pairs $(x,y)$, where $(x,y)$ and $(x',y')$ are equivalent when some $z$ satisfies $x+y'+z=x'+y+z$. Equality of $(M,0)$ and $(N,0)$ therefore says that $M+T=N+T$ for some $T$. This is exactly the asserted stable isomorphism.
\end{proof}

Cancellation would strengthen the conclusion to $M\cong N$. Roggenkamp's local cancellation results and Jacobinski's global cancellation theorem address when such a strengthening is valid. Their hypotheses are substantial. In particular, restrictions involving totally definite quaternion components are relevant precisely where quaternion orders also arise in supersingular elliptic-curve arithmetic. Removing these restrictions would erase part of the arithmetic content.\cite[Chapters VI--VIII]{v2Roggenkamp1970II}

\begin{example}[An invariant may deliberately forget torsion]
Let $G_0(\mathbb Z)$ use all finitely generated abelian groups. The exact sequence
\[
0\longrightarrow\mathbb Z\xrightarrow{\,n\,}\mathbb Z
 \longrightarrow\mathbb Z/n\mathbb Z\longrightarrow0
\]
shows $[\mathbb Z/n\mathbb Z]=0$. Every finitely generated abelian group is a sum of a free group and a finite group, so rank induces $G_0(\mathbb Z)\cong\mathbb Z$. By contrast, the Grothendieck group of the category of finite abelian groups has one independent generator for each prime, represented by $\mathbb Z/p\mathbb Z$. The category and the permitted exact sequences determine what the invariant remembers. This example is also among the exercises accompanying VIII.1 in the supplied volume.\cite{v2Roggenkamp1970II}
\end{example}

\section{Euler classes of complexes}

Let $C^\bullet$ be a bounded cochain complex in an abelian category:
\[
\cdots\longrightarrow C^{i-1}\xrightarrow{d^{i-1}}C^i
 \xrightarrow{d^i}C^{i+1}\longrightarrow\cdots,
\qquad d^i d^{i-1}=0.
\]
Write $Z^i=\ker d^i$, $B^i=\operatorname{im}d^{i-1}$, and $H^i=Z^i/B^i$. In a Grothendieck group containing these objects define
\[
\chi(C^\bullet)=\sum_i(-1)^i[C^i].
\]

\begin{theorem}[Algebraic Euler identity]
Suppose all terms, cycles, boundaries and cohomology objects belong to the exact category in question, and both sequences $0\to Z^i\to C^i\to B^{i+1}\to0$ and $0\to B^i\to Z^i\to H^i\to0$ are admissibly exact in its specified exact structure. Then
\[
\sum_i(-1)^i[C^i]=\sum_i(-1)^i[H^i(C^\bullet)].
\]
\end{theorem}
\begin{proof}
The two exact sequences
\[
0\to Z^i\to C^i\to B^{i+1}\to0,
\qquad
0\to B^i\to Z^i\to H^i\to0
\]
give $[C^i]=[B^i]+[H^i]+[B^{i+1}]$. Multiply by $(-1)^i$ and sum. Each boundary term occurs twice with opposite signs. Boundedness makes the sum finite and eliminates endpoint terms after inserting zero objects outside the given range.
\end{proof}

The hypothesis about the category is essential. A complex of lattices can have torsion cohomology, so its cohomology terms need not lie in $G_0^{\mathrm{lat}}$. The identity always makes sense in the Grothendieck group of all finitely generated modules when the ring is noetherian, but this is a different group. Similarly, a complex of projectives need not have projective cohomology, so one cannot insert its cohomology classes directly into $K_0(\Lambda)$ without a suitable resolution construction.

\begin{proposition}[Bounded acyclic projective complexes]
A bounded acyclic complex of projective modules is contractible. Its Euler class in $K_0(\Lambda)$ is zero.
\end{proposition}
\begin{proof}
Let $n$ be the largest degree with a nonzero term. Acyclicity makes $d^{n-1}:C^{n-1}\to C^n$ surjective. Since $C^n$ is projective, choose a section. Consequently $C^{n-1}\cong Z^{n-1}\oplus C^n$, and the last two terms contain a direct summand equal to the two-term identity complex $C^n\xrightarrow{1}C^n$. Its complement is a shorter acyclic complex of projectives: $Z^{n-1}$ is projective because it is a direct summand of $C^{n-1}$. Induction decomposes the entire complex into identity complexes. Each has an explicit contracting homotopy and contributes $[P]-[P]=0$ to the Euler class.
\end{proof}

For a bounded complex of projectives, this proves that the Euler class is invariant under quasi-isomorphism: the mapping cone of a quasi-isomorphism is a bounded acyclic projective complex, and its Euler class is the difference of the two Euler classes. This is the elementary core of the language of perfect complexes. It is enough for several useful arithmetic comparisons without claiming that every lattice has a finite projective resolution.

\begin{example}[A rationally exact integral complex]
For a prime $p$, consider
\[
0\longrightarrow\mathbb Z\xrightarrow{A}\mathbb Z^2
 \xrightarrow{B}\mathbb Z\longrightarrow0,
\quad A(t)=(pt,0),\quad B(u,v)=v.
\]
Then $BA=0$, $\ker A=0$, $\operatorname{coker}B=0$, and the middle cohomology is $\mathbb Z/p\mathbb Z$. Tensoring with $\mathbb Q$ makes the complex exact. Thus a generic-rank calculation over the fraction field detects none of its torsion defect. Modulo $p$, the first map becomes zero and the two nonzero cohomology groups each have dimension one. The alternating dimension remains zero throughout. Exactness, torsion and Euler characteristic are three different measurements of this example.
\end{example}

\section{Schanuel's lemma as a comparison mechanism}

\begin{theorem}[Schanuel's lemma]
Suppose
\[
0\to K\to P\xrightarrow{f}M\to0,
\qquad
0\to K'\to P'\xrightarrow{f'}M\to0
\]
are exact and $P,P'$ are projective. Then
\[
K\oplus P'\cong K'\oplus P.
\]
\end{theorem}
\begin{proof}
Form the pullback
\[
Q=\{(p,p')\in P\oplus P':f(p)=f'(p')\}.
\]
The projection $Q\to P'$ is surjective because $f$ is surjective, and its kernel is $K$. Since $P'$ is projective, $Q\cong K\oplus P'$. Interchanging the two presentations gives $Q\cong K'\oplus P$. Composing these isomorphisms proves the assertion.
\end{proof}

This lemma explains why alternating presentation classes are well defined even though the presentation itself is not unique. Roggenkamp uses it in VIII.3.1 to compare the Grothendieck group of lattices with that of all finitely generated modules. The point is not that a torsion module becomes a lattice; it is represented by a difference of two lattice classes.\cite[Theorem VIII.3.1]{v2Roggenkamp1970II}

\begin{proposition}[A one-step lattice presentation]
Every finitely generated left $\Lambda$-module $T$ admits an exact sequence
\[
0\to M\to P\to T\to0
\]
with $P$ finite free over $\Lambda$ and $M$ a $\Lambda$-lattice.
\end{proposition}
\begin{proof}
Choose finitely many generators of $T$ over $\Lambda$ to obtain a surjection $P=\Lambda^n\to T$. The ring $\Lambda$ is noetherian on either side because it is module-finite over the central noetherian ring $R$, so its kernel $M$ is finitely generated over $R$. It is an $R$-submodule of the torsion-free module $P$, hence is torsion-free. Therefore $M$ is a lattice.
\end{proof}

In particular $[P]-[M]$ is independent of the chosen one-step presentation by Schanuel's lemma. Additivity on exact sequences follows by the standard pullback construction of compatible presentations, yielding the comparison used in the book. This is a genuine arithmetic use of complexes and stable classes. It does not supply a short algorithm for finding module isomorphisms, since the existence of a presentation comparison is separate from the complexity of constructing it.

\section{What hereditary and Bass hypotheses change}

A left hereditary ring is one in which every left ideal is projective; equivalently every submodule of a projective left module is projective. For orders, this sharply simplifies lattice resolutions. A lattice can be embedded into a finite direct sum of copies of the ambient semisimple algebra and, after clearing denominators, into a finite free $\Lambda$-module. Under a hereditary hypothesis it is therefore projective. Here the embedding must be justified from the semisimplicity of the generic fiber; torsion-freeness alone over an arbitrary ring would not suffice. Roggenkamp's Chapter IX treats the structure of hereditary orders and explains why their local and global classification still requires care.\cite[Chapter IX, \S2]{v2Roggenkamp1970II}

\begin{proposition}[Embedding a lattice into a free module]
Let $\Lambda$ be an order in a finite-dimensional semisimple $K$-algebra $A$. Every $\Lambda$-lattice $M$ embeds into $\Lambda^n$ for some $n$.
\end{proposition}
\begin{proof}
The finite-dimensional $A$-module $K\otimes_R M$ is a direct summand of $A^n$ for a sufficiently large $n$, because $A$ is semisimple and its regular representation contains every simple module. Thus there is an injective $A$-linear map $j:K\otimes_RM\to A^n$. Choose finitely many $R$-generators $m_1,\ldots,m_\ell$ of $M$. Since $K\Lambda^n=A^n$, each $j(m_i)$ becomes an element of $\Lambda^n$ after multiplication by a nonzero denominator in $R$. Taking the product of these denominators gives a common nonzero $r\in R$ for which $rj(M)\subseteq\Lambda^n$. Multiplication by $r$ does not destroy injectivity in the $K$-vector space $A^n$. Restricting $rj$ to $M$ gives the desired embedding.
\end{proof}

\begin{corollary}
If $\Lambda$ is hereditary, every $\Lambda$-lattice is projective. Consequently every short exact sequence of lattices splits, and $K_0(\Lambda)\cong G_0^{\mathrm{lat}}(\Lambda)$.
\end{corollary}
\begin{proof}
The embedding proposition and the submodule characterization of hereditary rings make each lattice projective. A short exact sequence splits because its quotient is projective. Hence the two categories have the same objects and the same additive relations.
\end{proof}

Bass orders occupy a broader setting. In the standard order-theoretic terminology, an order is Gorenstein when its $R$-dual is projective on the appropriate opposite side, and it is Bass when every overorder is Gorenstein. The duality, finiteness and local hypotheses must be specified when using this terminology outside orders. Roggenkamp's 1971 paper computes Grothendieck groups of Bass orders over complete discrete valuation rings with finite residue field and distinguishes an exceptional local type. The paper also explains why a local surjectivity statement from maximal overorders does not simply persist globally.\cite{v2RoggenkampBass1971}

The useful conclusion here is methodological. A structural class can make a module problem more tractable and can expose stable invariants. One still has to determine whether the particular cryptographic order belongs to that class, which invariants survive a reduction, and what information an algorithm must recover. A theorem about finite representation type is not a theorem asserting polynomial-time recovery of a cryptographic secret.

\section{A Gram-matrix bridge already present in the book}

There is a second direct connection to the finite Hodge calculations below. In the proof of VIII.3.5, Roggenkamp compares a decomposition matrix and a Cartan matrix by means of a positively weighted Gram matrix. On printed page 139, the argument explicitly uses the equality between the rank of a real matrix and the rank of its Gram matrix. This is a concrete positivity argument inside the supplied arithmetic source. It is not the analytic index theorem, but it is the same elementary mechanism that identifies harmonic vectors with simultaneous differential kernels.\cite[Theorem VIII.3.5 and Exercise VIII.3.5]{v2Roggenkamp1970II}

\begin{proposition}[Weighted Gram kernel and rank]
Let $X$ be an $m\times n$ real matrix, and let $W=\operatorname{diag}(w_1,\ldots,w_m)$ with every $w_i>0$. Then
\[
\ker(X^TWX)=\ker X,
\qquad
\operatorname{rank}(X^TWX)=\operatorname{rank}X.
\]
Moreover $X^TWX$ is positive semidefinite.
\end{proposition}
\begin{proof}
For every column vector $v\in\mathbb R^n$,
\[
v^TX^TWXv=\sum_{i=1}^m w_i(Xv)_i^2\geq0.
\]
Symmetry is immediate from the diagonal form of $W$, proving positive semidefiniteness. If $Xv=0$, then $X^TWXv=0$. Conversely, if $X^TWXv=0$, the displayed sum is zero. Each summand is nonnegative and each weight is strictly positive, so every $(Xv)_i$ is zero. This proves the equality of kernels. Both matrices define linear maps with the same $n$-dimensional domain, so rank--nullity gives equality of ranks.
\end{proof}

The positive weights matter. If $W$ is merely positive semidefinite, nonzero coordinates of $Xv$ in a zero-weight direction can disappear. If $W$ is indefinite, cancellation can occur between positive and negative contributions. For example, $X=(1,1)^T$ and $W=\operatorname{diag}(1,-1)$ give $X^TWX=0$ although $X$ has rank one. Similarly, over a finite field a nonzero vector can have zero sum of squares. Thus the real positivity step in the book is not a formal identity that can be transferred unchanged to arbitrary coefficient fields.

For a real two-step complex with chosen inner products, place $A^*$ and $B$ in a single vertically stacked matrix $X$. Its unweighted Gram matrix is $AA^*+B^*B$, the middle Hodge Laplacian. The proposition then identifies its kernel with $\ker A^*\cap\ker B$. This gives a precise arithmetic-to-Hodge connection at the level of a matrix lemma. Its proof requires no claim that the arithmetic category and the differential-geometric category are equivalent. It also supplies no scale-independent estimate: equality of kernels and ranks does not determine the smallest positive singular value.

\section{From projective modules to topological bundles}

For a compact Hausdorff space $X$, let $C(X,\mathbb C)$ be the ring of continuous complex-valued functions. Swan's theorem identifies finite-rank complex vector bundles on $X$ with finitely generated projective modules over $C(X,\mathbb C)$ through the section functor. It induces
\[
K_0\bigl(C(X,\mathbb C)\bigr)\cong K^0(X).
\]
This is a precise comparison theorem, not an equality between arbitrary algebraic and topological constructions.\cite{v2Swan1962}

Its idempotent mechanism is worth deriving. A finitely generated projective module $P$ is the image of an idempotent matrix $e\in M_n(C(X,\mathbb C))$. At each $x\in X$, the matrix $e(x)$ satisfies $e(x)^2=e(x)$. Its image is a vector subspace of $\mathbb C^n$. The rank is locally constant: the eigenvalues of an idempotent are zero and one, and its trace is an integer-valued continuous function. Around a point where an appropriate minor is nonzero, its image admits a continuous frame. These image spaces therefore form a vector bundle. Conversely, a vector bundle on a compact space admits a complementary bundle making it a summand of a trivial bundle; the resulting projection gives an idempotent. This describes how the two kinds of object meet.

For an arithmetic order $\Lambda$, one cannot silently replace $\Lambda$ by $C(X,\mathbb C)$. One needs an actual representation or scalar-extension construction. For example, a unital homomorphism
\[
\rho:\Lambda^{\mathrm{op}}\longrightarrow M_n(C(X,\mathbb C))
\]
specifies a right $\Lambda$-action on the trivial rank-$n$ bundle. Tensoring it with a finitely generated projective left $\Lambda$-module then produces a vector bundle, because the latter module is an idempotent summand of a finite free module. This yields a homomorphism on $K_0$. Its definition depends on $\rho$; no canonical map is supplied merely by knowing that both domains involve K-theory.

\begin{remark}[A left--right convention]
A homomorphism into endomorphisms gives a left action by default. To use it on the right, one must provide a homomorphism from $\Lambda^{\mathrm{op}}$, or equivalently specify a right action directly. Throughout this scalar-extension construction the data are a bundle $E$ with a right $\Lambda$-action and a left projective $\Lambda$-module $P$, and the output is $E\otimes_\Lambda P$. This convention avoids reversing the multiplication in a noncommutative order.
\end{remark}

There is also an important degree-one distinction. Algebraic $K_1(\Lambda)$ is formed from the stable general linear group modulo elementary matrices. Topological $K_1$ uses a homotopy construction. For the one-point space,
\[
K_1^{\mathrm{alg}}(\mathbb C)\cong\mathbb C^\times,
\qquad K_1^{\mathrm{top}}(\mathbb C)=0.
\]
The first equality follows from the determinant and elementary elimination; the second follows because complex invertible matrices are path connected. The agreement in degree zero under Swan's theorem therefore does not justify identifying all higher algebraic and topological K-groups. Likewise the compact space is recovered as the character or maximal-ideal space of its continuous-function algebra with the appropriate topology; the full algebraic prime spectrum is a different object. These distinctions refine the historical discussion on printed page 114 of the supplied volume.\cite{v2Roggenkamp1970II,v2Swan1962}

\chapter{Varieties of complexes, exactness, and deformation}
\label{ch:varieties-complexes}

\section{A parameter space with an equation}

Fix finite-dimensional vector spaces $V_0,V_1,V_2$ over a field $k$, with dimensions $a,b,c$. A pair of maps
\[
V_0\xrightarrow{A}V_1\xrightarrow{B}V_2
\]
is a complex precisely when $BA=0$. After choosing bases, its parameter scheme is
\[
\mathcal C_{a,b,c}
=\operatorname{Spec}
 k[a_{ij},b_{\ell i}]/\bigl((BA)_{\ell j}:1\leq\ell\leq c,
                                      1\leq j\leq a\bigr).
\]
The equations are quadratic. They express compatibility of successive maps, but they do not assert exactness. Exactness is an additional condition on images and kernels.

Buchsbaum and Eisenbud developed structure and exactness results for finite free complexes. De Concini and Strickland studied the algebraic geometry of varieties of complexes; the later \emph{Hodge algebras} exposition by De Concini, Eisenbud and Procesi gives a direct account in Section 16. It is therefore preferable to identify the precise result being used instead of attributing the entire geometric theory to one name. The calculations below are direct derivations for two successive maps and require neither the full straightening-law machinery nor a classification of all singularities.\cite{v2BE1973,v2DS1981,v2DEP1982}

\begin{proposition}[Rank constraint and cohomology]
If $BA=0$, $r=\operatorname{rank}A$, and $s=\operatorname{rank}B$, then
\[
r+s\leq b.
\]
The cohomology dimensions of the three-term complex are
\[
h^0=a-r,\qquad h^1=b-r-s,\qquad h^2=c-s.
\]
In particular the complex is exact at $V_1$ if and only if $r+s=b$.
\end{proposition}
\begin{proof}
The equation $BA=0$ implies $\operatorname{im}A\subseteq\ker B$. Rank--nullity gives $\dim\ker B=b-s$, so $r\leq b-s$. The formula for $h^1$ is the dimension of the quotient $\ker B/\operatorname{im}A$. The endpoint formulas are rank--nullity for $A$ and the dimension of the cokernel of $B$. Equality of the middle image and kernel follows from inclusion and equality of their dimensions.
\end{proof}

\begin{corollary}[Euler characteristic]
For every pair in $\mathcal C_{a,b,c}$,
\[
h^0-h^1+h^2=a-b+c.
\]
\end{corollary}
\begin{proof}
Substitution cancels $r$ and $s$.
\end{proof}

This invariant is constant even when cohomology jumps. It is the numerical version of the Euler identity in the preceding chapter. Its constancy is useful, but it is exactly why it cannot by itself detect a jump or an estimate on small singular values.

\section{Classification under changes of bases}

The group
\[
G=\operatorname{GL}(V_0)\times\operatorname{GL}(V_1)
   \times\operatorname{GL}(V_2)
\]
acts by
\[
(g_0,g_1,g_2)\cdot(A,B)
=(g_1Ag_0^{-1},\ g_2Bg_1^{-1}).
\]
This action preserves $BA=0$ because the composite transforms to $g_2BAg_0^{-1}$. It represents isomorphism of complexes with the graded spaces fixed up to independent bases.

\begin{theorem}[Normal form over a field]
Two points of $\mathcal C_{a,b,c}$ lie in the same $G$-orbit if and only if they have the same pair of ranks $(r,s)$.
\end{theorem}
\begin{proof}
Ranks are unchanged by the action. Conversely, choose a decomposition
\[
V_0=H_0\oplus U,\qquad
V_1=I\oplus H_1\oplus W,\qquad
V_2=J\oplus H_2,
\]
where $H_0=\ker A$, $I=\operatorname{im}A$, $\ker B=I\oplus H_1$, and $J=\operatorname{im}B$. Choose $U$ complementary to $H_0$ and $W$ complementary to $\ker B$. The restrictions $A:U\to I$ and $B:W\to J$ are isomorphisms. Choose bases of $U$ and $W$, transport them to $I$ and $J$, and choose arbitrary bases of $H_0,H_1,H_2$. In these bases the only nonzero blocks of $A$ and $B$ are identity matrices of sizes $r$ and $s$. The dimensions of all blocks are determined by $(a,b,c,r,s)$, so any two such pairs have the same normal form.
\end{proof}

This theorem depends strongly on the coefficient field. It fails over an order or over $\mathbb Z$ if ranks are the only invariants retained. The factor $p$ in the integral example above is invisible over $\mathbb Q$ and essential over $\mathbb Z$. A visual rank classification over $\mathbb R$ or $\mathbb C$ must therefore be presented as a model of the linear-algebraic portion of the theory, not as a complete classification of integral lattice complexes.

\section{Dimension of a rank stratum}

Let $\mathcal C^{r,s}_{a,b,c}$ be the locally closed set where $A$ has rank exactly $r$ and $B$ has rank exactly $s$. It is nonempty exactly when
\[
0\leq r\leq\min(a,b),\quad
0\leq s\leq\min(b,c),\quad r+s\leq b.
\]
The normal form proves sufficiency. The rank inequalities prove necessity.

\begin{theorem}[Dimension formula]
Over an algebraically closed field, every nonempty rank stratum is smooth and irreducible, and
\[
\dim\mathcal C^{r,s}_{a,b,c}
=r(a+b-r)+s(b-r+c-s).
\]
\end{theorem}
\begin{proof}
First choose $I=\operatorname{im}A\subseteq V_1$, a point of the Grassmannian $\operatorname{Gr}(r,V_1)$ of dimension $r(b-r)$. For this $I$, the choices of a surjection $V_0\to I$ form an open subset of $\operatorname{Hom}(V_0,I)$, of dimension $ar$. Thus the rank-$r$ locus for $A$ has dimension $r(a+b-r)$ and is smooth and irreducible.

For fixed $A$, the equation $BA=0$ says precisely that $B$ factors through $V_1/I$. A rank-$s$ map from this $(b-r)$-dimensional quotient to $V_2$ is specified by its image $J\in\operatorname{Gr}(s,V_2)$ and a surjection $V_1/I\to J$. The choices have dimension $s(c-s)+s(b-r)$. Locally on Grassmannian charts, the quotient spaces form a trivial vector bundle and these choices form a product with the rank-$s$ matrix locus. This gives smoothness and the displayed dimension. The Grassmannians and the open surjection loci are irreducible, and the bundle description preserves irreducibility.
\end{proof}

The formula is a count of algebraic parameters, not a complexity bound. A search space of large dimension can still be easy to parameterize, while a finite set can encode a hard search problem. For cryptographic use one must also fix the field size or integer height, the input encoding, and which maps or invariants are visible to an adversary.

\begin{example}[The first nontrivial quadric]
Take $(a,b,c)=(1,2,1)$. With $A=(x_1,x_2)^T$ and $B=(y_1,y_2)$, the sole equation is
\[
q=x_1y_1+x_2y_2=0.
\]
The $(1,1)$-stratum has dimension $1(1+2-1)+1(2-1+1-1)=3$. It is the locus where neither $x$ nor $y$ is zero. The $(1,0)$ and $(0,1)$ strata each have dimension two, and $(0,0)$ is the origin. Every point in the $(1,1)$-stratum is an exact three-term complex.
\end{example}

\section{First-order deformations and tangent equations}

Let $k[\varepsilon]/(\varepsilon^2)$ be the ring of dual numbers. A first-order perturbation of $(A,B)$ has the form
\[
A_\varepsilon=A+\varepsilon U,
\qquad B_\varepsilon=B+\varepsilon V.
\]
Multiplication gives
\[
B_\varepsilon A_\varepsilon
=BA+\varepsilon(BU+VA).
\]
Thus the Zariski tangent space at a point with $BA=0$ is
\[
T_{(A,B)}\mathcal C_{a,b,c}
=\{(U,V):BU+VA=0\}.
\]
This calculation is valid over any field. It is the derivative of the polynomial equation defining the scheme, so it applies to the scheme tangent space whether or not the point is smooth.

\begin{proposition}[Tangent-space dimension]
At a point of rank $(r,s)$,
\[
\dim T_{(A,B)}\mathcal C_{a,b,c}
=ab+bc-as-cr+rs.
\]
\end{proposition}
\begin{proof}
Consider the linear map
\[
L:\operatorname{Hom}(V_0,V_1)\oplus\operatorname{Hom}(V_1,V_2)
 \longrightarrow\operatorname{Hom}(V_0,V_2),
\quad L(U,V)=BU+VA.
\]
The maps $BU$ are exactly the maps $V_0\to\operatorname{im}B$. Their space has dimension $as$: every such map lifts through the surjection $B:V_1\to\operatorname{im}B$, since vector-space sequences split. The maps $VA$ are exactly the maps $V_0\to V_2$ that vanish on $\ker A$. Their space has dimension $cr$, since $A$ identifies $V_0/\ker A$ with $\operatorname{im}A$ and a map on that image extends to $V_1$.

The intersection consists of maps from $V_0/\ker A$ to $\operatorname{im}B$, so has dimension $rs$. Therefore $\operatorname{rank}L=as+cr-rs$. The tangent space is $\ker L$, and rank--nullity gives the result.
\end{proof}

\begin{corollary}[Transverse infinitesimal directions]
The difference between tangent-space dimension and rank-stratum dimension is
\[
(b-r-s)(a+c-r-s)=h^1(h^0+h^2).
\]
\end{corollary}
\begin{proof}
Subtract the preceding dimension formulas and factor; substitute the cohomology dimensions already proved.
\end{proof}

This difference measures directions that leave the fixed-rank stratum to first order. Its positivity alone does not prove that the ambient variety is singular: a lower-dimensional stratum can lie inside a smooth ambient locus. For the quadric $q=0$ with $(a,b,c)=(1,2,1)$, a point with $A\neq0$ and $B=0$ lies in a two-dimensional stratum but is smooth on the three-dimensional quadric.

\begin{proposition}[The singularity of the four-coordinate model]
Assume $\operatorname{char}k\neq2$. The hypersurface
\[
x_1y_1+x_2y_2=0\subseteq\mathbb A_k^4
\]
is smooth away from the origin and singular at the origin.
\end{proposition}
\begin{proof}
Its gradient is $(y_1,y_2,x_1,x_2)$. This vector is zero exactly at the origin. Away from the origin the hypersurface Jacobian has rank one, yielding a three-dimensional tangent space and smoothness. At the origin the tangent space is the entire four-dimensional ambient space, while the hypersurface has dimension three. Thus the origin is singular.
\end{proof}

The proof in fact extends to this particular polynomial in characteristic two using the same derivatives and the hypersurface criterion. The characteristic restriction is retained here because the orthogonal coordinate change used in the real visualization below is part of the model being illustrated.

\section{The endomorphism complex and the meaning of gauge}

There is a uniform way to package these first-order calculations. For a bounded graded vector space $C^\bullet$, define
\[
\operatorname{End}^m(C)
=\bigoplus_i\operatorname{Hom}(C^i,C^{i+m}).
\]
If $d$ is a differential of degree one, put
\[
\delta f=df-(-1)^mfd
\quad\text{for }f\in\operatorname{End}^m(C).
\]
Composition is the multiplication; the formula is a graded commutator.

\begin{proposition}[Endomorphism differential]
If $d^2=0$, then $\delta^2=0$. For homogeneous $f$ of degree $m$ and any homogeneous $g$,
\[
\delta(fg)=(\delta f)g+(-1)^m f(\delta g).
\]
\end{proposition}
\begin{proof}
For $f$ of degree $m$,
\begin{align*}
\delta^2f
&=d\bigl(df-(-1)^mfd\bigr)
 -(-1)^{m+1}\bigl(df-(-1)^mfd\bigr)d\\
&=d^2f-(-1)^mdfd+(-1)^mdfd-fd^2=0.
\end{align*}
For the second identity expand both sides. The two terms containing $fdg$ appear with opposite signs and cancel, leaving $dfg-(-1)^{m+\deg g}fgd$.
\end{proof}

A first-order change $d_\varepsilon=d+\varepsilon a$ satisfies $d_\varepsilon^2=0$ if and only if $\delta a=da+ad=0$. A degree-zero change of basis $g_\varepsilon=1+\varepsilon u$ has inverse $1-\varepsilon u$, and
\[
g_\varepsilon d g_\varepsilon^{-1}
=d+\varepsilon(ud-du)=d-\varepsilon\delta u.
\]
Consequently first-order deformations modulo infinitesimal changes of basis are described by $H^1(\operatorname{End}^\bullet(C),\delta)$.

The word ``gauge'' in this calculation means a change of bases in a graded algebraic object. In a differential-geometric setting, a gauge transformation is a bundle automorphism and can vary over a manifold. The formulas have a common deformation-theoretic pattern, but the analytic function spaces, global topology and equations of motion are additional data. We will derive the geometric version explicitly rather than treating the terminology as a proof of equivalence.

\begin{proposition}[Second-order obstruction]
Let $a\in\operatorname{End}^1(C)$ satisfy $\delta a=0$. A lift
\[
d_t=d+ta+t^2b\pmod{t^3}
\]
is a differential if and only if $\delta b=-a^2$. In particular the class of $a^2$ in $H^2(\operatorname{End}^\bullet(C),\delta)$ must vanish, and its vanishing is sufficient for such a second-order lift.
\end{proposition}
\begin{proof}
Expanding the square gives
\[
d_t^2=d^2+t(da+ad)+t^2(db+bd+a^2)\pmod{t^3}.
\]
The first two coefficients vanish by hypothesis. The third vanishes precisely when $\delta b=-a^2$. The graded Leibniz rule gives $\delta(a^2)=(\delta a)a-a(\delta a)=0$, so $a^2$ is a cocycle. Existence of $b$ is exactly the assertion that its cohomology class is zero.
\end{proof}

Over orders, extension and deformation groups can have torsion, and reduction modulo a prime can change their dimensions. Thus the field calculation should be regarded as one layer in a larger integral deformation problem. It suggests a concrete research workflow: construct an integral complex, compute its cohomology and endomorphism cohomology, compare their reductions, and only then investigate whether the deformation parameters have any relation to a cryptographic task.

\section{Why ranks alone fail over commutative rings}

For a matrix $M$ with entries in a commutative ring $S$, let $I_t(M)$ be the ideal generated by its $t\times t$ minors, with $I_0(M)=S$. Minors behave well under base change: applying a ring homomorphism to the matrix applies the homomorphism to every minor. The condition $I_t(M)=S$ is much stronger than saying that one $t\times t$ minor is nonzero. It says that locally everywhere some such minor is invertible.

\begin{lemma}[Splitting off a unit pivot]
Let $S$ be a commutative ring and let a differential in a complex of finite free $S$-modules have a unit matrix entry. After changes of bases the complex has a direct summand consisting of $S\xrightarrow{1}S$ in the two affected degrees.
\end{lemma}
\begin{proof}
Permute rows and columns to place the unit in the upper-left entry and scale its row to make it one. Elementary operations clear the rest of its column and row. The differential then has block form $\operatorname{diag}(1,M')$. The identity $d^id^{i-1}=0$ forces the component of the preceding map into this pivot coordinate to be zero. Likewise $d^{i+1}d^i=0$ forces the following map to vanish on its image coordinate. The two coordinates therefore split off as an identity complex.
\end{proof}

This lemma is the ring-theoretic remnant of Gaussian elimination. A nonzero scalar is invertible over a field, so one can keep eliminating. Over a ring, a nonzero scalar such as $p\in\mathbb Z$ need not be invertible, and its residual arithmetic becomes part of the homology.

For a noetherian local ring $(S,\mathfrak m)$ and a proper ideal $I$, its grade is the maximum length of an $S$-regular sequence contained in $I$. A sequence $x_1,\dots,x_t$ is regular if each $x_j$ is a nonzerodivisor on $S/(x_1,\dots,x_{j-1})$ and the final quotient is nonzero. We set $\operatorname{grade}S=+\infty$ for the unit ideal. Over a general noetherian ring the same concept can be expressed by localized depth or Ext; the local formulation is sufficient here.

\begin{theorem}[Buchsbaum--Eisenbud criterion; external theorem]
Let $S$ be a commutative noetherian local ring and let
\[
0\to F_n\xrightarrow{d_n}F_{n-1}\to\cdots
 \xrightarrow{d_2}F_1\xrightarrow{d_1}F_0
\]
be a complex of finite free modules. Put
\[
r_i=\sum_{j=i}^n(-1)^{j-i}\operatorname{rank}F_j.
\]
Using the usual determinantal rank and the convention for the unit ideal, the complex is acyclic in positive degrees if and only if
\[
\operatorname{rank}d_i=r_i,
\qquad \operatorname{grade}I_{r_i}(d_i)\geq i
\quad(1\leq i\leq n).
\]
Here acyclic in positive degrees leaves $\operatorname{coker}d_1$ unconstrained. Zero free summands can be removed, or handled with the standard zero-rank conventions.\cite{v2BE1973}
\end{theorem}

The full theorem is an external result. The manuscript does not claim that its proof is a few lines of rank--nullity, and the accompanying formalization should not replace it by an unproved axiom while advertising a verified proof. Its two ingredients have different roles: rank conditions control generic compatibility; grade conditions prevent homology from surviving on loci of insufficient codimension or depth.

The integral example $A=(p,0)^T$, $B=(0,1)$ illustrates the second ingredient after localizing $\mathbb Z$ at $(p)$. The ideal of rank-one minors of $A$ is $(p)$, with grade one. The leftmost map is in homological degree two, so the required grade is at least two. The criterion correctly rejects exactness. Over $\mathbb Q$ the same ideal becomes the unit ideal, and the obstruction disappears. The change is arithmetic, not a disagreement between two rank computations.

\section{A complete exactness proof: the two-generator Koszul complex}

Let $S$ be a commutative ring and let $x,y\in S$. Define
\[
d_2:S\to S^2,\quad d_2(c)=(-yc,xc),
\qquad
 d_1:S^2\to S,\quad d_1(a,b)=xa+yb.
\]
Then $d_1d_2=0$ by commutativity.

\begin{theorem}[Regular two-generator Koszul complex]
Suppose $x$ is a nonzerodivisor in $S$ and the class of $y$ is a nonzerodivisor in $S/(x)$. Then
\[
0\to S\xrightarrow{d_2}S^2\xrightarrow{d_1}S
 \to S/(x,y)\to0
\]
is exact.
\end{theorem}
\begin{proof}
If $d_2(c)=0$, then $xc=0$, hence $c=0$. Thus $d_2$ is injective. Its image is contained in $\ker d_1$ by the complex identity. Conversely, suppose $xa+yb=0$. Modulo $(x)$ this becomes $\overline y\,\overline b=0$. The hypothesis gives $b=xc$ for some $c\in S$. Substitution yields $x(a+yc)=0$, hence $a=-yc$. Therefore $(a,b)=d_2(c)$, proving exactness at $S^2$. Finally the image of $d_1$ is exactly the ideal $(x,y)$, which is the kernel of the quotient map to $S/(x,y)$.
\end{proof}

The proof shows precisely where regularity is used. Merely having two nonzero elements is insufficient. If $S=k[t]/(t^2)$ and $x=t$, multiplication by $x$ already has a nonzero kernel; no generic intuition can remove that kernel. If $S=k[x,y]$, the variables form a regular sequence and the preceding resolution is exact. This example connects the local algebra of free resolutions with the symbol calculations of elliptic complexes: both use exterior-algebra differentials, but regularity in a ring and nonzero covectors on a manifold are different hypotheses.

\section{Base change and the appearance of new homology}

Let $S\to T$ be a ring homomorphism. Tensoring a complex with $T$ preserves the equation $d^2=0$, because composition and tensor product commute. It need not preserve exactness.

\begin{proposition}[Flat base change]
If $T$ is flat as an $S$-module, tensoring an exact complex of $S$-modules with $T$ preserves exactness. For any complex, the natural comparison identifies its cohomology after tensoring with $H^i(C)\otimes_ST$.
\end{proposition}
\begin{proof}
Flatness preserves the two short exact sequences involving cycles, boundaries and cohomology used in the Euler identity. Their tensor products identify the new cycles and boundaries with the tensors of the old ones. Taking their quotient gives the cohomology comparison. When the original cohomology vanishes, the new one vanishes as well.
\end{proof}

Localization is flat, so passing from an integral domain to its fraction field preserves exactness. The converse fails because localization can kill torsion. Reduction modulo a prime is usually not flat. For example, tensor the exact sequence
\[
0\to\mathbb Z\xrightarrow{\,p\,}\mathbb Z
 \to\mathbb Z/p\mathbb Z\to0
\]
with $\mathbb F_p$. The multiplication map becomes zero and loses injectivity. This loss is measured by a Tor term, not by an error in the tensor operation.

For an arithmetic application, one should record at least three coefficient regimes: the order itself, its localizations or completions, and its residue fields. The same written matrix can define complexes with different homology in these regimes. A Lean theorem over a field verifies the field statement; it does not automatically establish the corresponding integral statement. A well-designed theorem interface makes the coefficient assumptions explicit enough that this distinction is enforced by type checking.

\section{The real four-dimensional quadric as an exact visual model}

Return to $q(x,y)=x\cdot y$ with $x,y\in\mathbb R^2$. Define
\[
u=\frac{x+y}{\sqrt2},\qquad
v=\frac{x-y}{\sqrt2}.
\]
Direct expansion gives
\[
q(x,y)=\frac{\|u\|^2-\|v\|^2}{2},
\qquad
\|x\|^2+\|y\|^2=\|u\|^2+\|v\|^2.
\]
The variety $q=0$ is therefore the null cone of a quadratic form of signature $(2,2)$. Its intersection with a sphere of radius $\rho>0$ is described by
\[
\|u\|=\|v\|=\frac{\rho}{\sqrt2}.
\]
This link is a product of two circles. An exact parameterization is
\[
u=\frac{\rho}{\sqrt2}(\cos\alpha,\sin\alpha),
\quad
v=\frac{\rho}{\sqrt2}(\cos\beta,\sin\beta),
\quad
x=\frac{u+v}{\sqrt2},\quad y=\frac{u-v}{\sqrt2}.
\]
It supplies a genuine four-coordinate object whose three-dimensional projections can be rotated interactively. The displayed projection can conceal intersections or identify points that are distinct in four dimensions; the underlying coordinates and the residual $q$ should therefore remain visible in the interface.

There are several useful controls with an exact mathematical interpretation. Changing $\rho$ moves along the cone and reaches the singular point at $\rho=0$. Changing $\alpha-\beta$ changes the relative sizes of $x$ and $y$; some values make one vanish and move to an endpoint rank stratum. A four-dimensional orthogonal rotation changes only the viewing coordinates if the original coordinates are retained for evaluating $q$. An arbitrary ambient rotation does not necessarily preserve the fixed expression $x\cdot y$, so the software must distinguish a change of view from a change of the mathematical model.

This example does not illustrate every variety of complexes. It illustrates the smallest nontrivial three-term case in four real coordinates. Its value is that every visible feature has an exact equation, a rank interpretation, and a verifiable relation to cohomology. It can be used to test the visualization code with known cases before any higher-dimensional object is drawn.

\chapter{Elliptic complexes, Hodge theory, and gauge deformations}
\label{ch:elliptic-complexes}

\section{Ellipticity is a condition on symbols}

Let $X$ be a closed smooth manifold: throughout this chapter ``closed'' means compact and without boundary. Let $E_i\to X$ be smooth complex vector bundles, and let
\[
D_i:\Gamma(E_i)\longrightarrow\Gamma(E_{i+1})
\]
be first-order differential operators satisfying $D_{i+1}D_i=0$. The collection is a differential complex. For a covector $\xi\in T_x^*X$, the principal symbol $\sigma(D_i)(x,\xi)$ records the highest-derivative part of $D_i$ in direction $\xi$. One convention is
\[
\sigma(D)(x,df_x)s_x=\sqrt{-1}\bigl(D(fs)_x-f(x)D(s)_x\bigr).
\]
Omitting the factor of $\sqrt{-1}$ gives another common convention; exactness of the symbol sequence is unchanged. The displayed analytic convention makes the Laplacian symbol positive.

\begin{definition}[Elliptic complex]
The differential complex is elliptic if, for every $(x,\xi)$ with $\xi\neq0$, the finite-dimensional sequence
\[
0\to E_{0,x}\xrightarrow{\sigma(D_0)(x,\xi)}E_{1,x}
 \to\cdots\to E_{m,x}\to0
\]
is exact.
\end{definition}

Thus an elliptic complex produces, at every nonzero cotangent direction, a point in an exact rank locus of a variety of complexes. This is a direct connection to the preceding chapter. The converse needs additional geometric information: a family of exact matrices is not automatically the symbol of differential operators, and even an admissible symbol need not have a chosen operator realization satisfying $D_{i+1}D_i=0$.

Ellipticity also differs from an elliptic curve. A smooth projective genus-one curve with a chosen point is an algebraic object. An elliptic differential operator is defined by the invertibility or exactness of its principal symbol away from the zero section. Elliptic curves carry useful elliptic complexes, such as their Dolbeault complex, but the two uses of the word do not identify their theories.

\section{Exterior algebra gives a complete symbol calculation}

Let $V$ be a finite-dimensional real or complex vector space. For $\xi\in V^*$ let $\epsilon(\xi)$ be exterior multiplication by $\xi$ on $\bigwedge^\bullet V^*$. For $v\in V$, let $\iota_v$ be contraction.

\begin{lemma}[Wedge--contraction identity]
For every $\alpha\in\bigwedge^\bullet V^*$,
\[
\iota_v\bigl(\xi\wedge\alpha\bigr)
+\xi\wedge\iota_v\alpha=\xi(v)\alpha.
\]
\end{lemma}
\begin{proof}
It suffices by linearity to take $\alpha=\eta_1\wedge\cdots\wedge\eta_p$. The definition of contraction is
\[
\iota_v\alpha=\sum_{j=1}^p(-1)^{j-1}\eta_j(v)
 \eta_1\wedge\cdots\widehat{\eta_j}\cdots\wedge\eta_p.
\]
Apply this definition to $\xi\wedge\alpha$. The contribution from the first factor is $\xi(v)\alpha$; every other contribution is the negative of the corresponding term in $\xi\wedge\iota_v\alpha$. Adding the two expressions cancels those remaining terms.
\end{proof}

\begin{theorem}[Exactness of the nonzero-covector Koszul complex]
If $\xi\neq0$, the complex
\[
0\to\bigwedge^0V^*\xrightarrow{\epsilon(\xi)}\bigwedge^1V^*
 \xrightarrow{\epsilon(\xi)}\cdots
 \xrightarrow{\epsilon(\xi)}\bigwedge^{\dim V}V^*\to0
\]
is contractible and hence exact.
\end{theorem}
\begin{proof}
Choose $v$ with $\xi(v)=1$. The wedge--contraction identity says that $\iota_v$ is a contracting homotopy. More explicitly, if $\xi\wedge\alpha=0$, then $\alpha=\xi\wedge\iota_v\alpha$, so each cycle is a boundary. The endpoint statements are included in the same identity.
\end{proof}

For the de Rham differential $d$, the symbol is $\sqrt{-1}\epsilon(\xi)$; the nonzero scalar does not change exactness. The theorem proves ellipticity of the de Rham complex. It is a pointwise proof and does not require a global coordinate system on $X$. The construction shows exactly how the algebraic condition $d^2=0$ and the analytic symbol condition work together: the former makes cohomology meaningful, while the latter supplies local high-frequency control.

With a Riemannian metric, contraction by the metric dual $\xi^\sharp$ is the adjoint of wedge multiplication by $\xi$. The same identity becomes
\[
\epsilon(\xi)\epsilon(\xi)^*
+\epsilon(\xi)^*\epsilon(\xi)
=\|\xi\|^2\operatorname{id}.
\]
This is the symbol of the Hodge Laplacian. Its positivity for $\xi\neq0$ is an ellipticity statement. It is not yet a quantitative lower bound on the nonzero spectrum of the global Laplacian.

\section{Finite-dimensional Hodge theory with a full proof}

Before invoking analytic elliptic theory, one can prove its finite-dimensional model completely. Give a finite cochain complex $C^\bullet$ Hermitian inner products. Let $(d^i)^*$ denote the adjoint, and define
\[
\Delta_i=(d^i)^*d^i+d^{i-1}(d^{i-1})^*.
\]
A vector is harmonic if it belongs to $\ker\Delta_i$.

\begin{theorem}[Finite-dimensional Hodge decomposition]
For each degree $i$,
\[
C^i=\operatorname{im}d^{i-1}\ \mathbin{\perp}\
\ker\Delta_i\ \mathbin{\perp}\ \operatorname{im}(d^i)^*,
\]
and the map from a harmonic vector to its cohomology class is an isomorphism
\[
\ker\Delta_i\cong H^i(C^\bullet).
\]
\end{theorem}
\begin{proof}
For $v\in C^i$,
\[
\langle\Delta_i v,v\rangle
=\|d^iv\|^2+\|(d^{i-1})^*v\|^2.
\]
Thus $\ker\Delta_i=\ker d^i\cap\ker(d^{i-1})^*$. The two image spaces are orthogonal because
\[
\langle d^{i-1}u,(d^i)^*w\rangle
=\langle d^id^{i-1}u,w\rangle=0.
\]
Their common orthogonal complement is exactly $\ker(d^{i-1})^*\cap\ker d^i$, proving the orthogonal direct-sum decomposition.

If $v$ is a cycle, write it in that decomposition as $v=d^{i-1}u+h+(d^i)^*w$. The last term is itself a cycle because the other two terms are. Taking its inner product with itself gives
\[
\|(d^i)^*w\|^2=\langle w,d^i(d^i)^*w\rangle=0.
\]
Therefore $v$ differs from the harmonic vector $h$ by a boundary. This proves surjectivity onto cohomology. A harmonic boundary has zero inner product with itself because harmonic vectors are orthogonal to boundaries, so it is zero. This proves injectivity.
\end{proof}

There is a spectral consequence with equally clear limits. On the orthogonal complement of its kernel, a fixed finite-dimensional $\Delta_i$ has a positive smallest eigenvalue whenever that complement is nonzero. Positivity follows from finite dimensionality and positive definiteness on that subspace. A family of such matrices need not have a uniform positive lower bound.

\begin{example}[A fixed index and a variable spectral gap]
In the real four-coordinate model, suppose $x\cdot y=0$. The three Hodge Laplacians are
\[
\Delta_0=\|x\|^2,
\qquad \Delta_1=xx^T+yy^T,
\qquad \Delta_2=\|y\|^2.
\]
When $x$ and $y$ are both nonzero, they are orthogonal and span $\mathbb R^2$. The eigenvalues of $\Delta_1$ are $\|x\|^2$ and $\|y\|^2$, because applying it to $x$ gives $\|x\|^2x$ and similarly for $y$. The complex is exact and its smallest eigenvalue on the total graded space is
\[
\min\{\|x\|^2,\|y\|^2\}.
\]
Taking $x=(t,0)$ and $y=(0,1)$ with $t>0$ preserves exactness and Euler characteristic zero, while the gap tends to zero as $t\to0$. At $t=0$ cohomology appears. An index can remain fixed through this entire process.
\end{example}

\section{The analytic Hodge theorem and its input}

For an elliptic complex on a closed manifold, choose Hermitian bundle metrics and a positive volume form. These define formal adjoints $D_i^*$. The Laplacians
\[
\Delta_i=D_i^*D_i+D_{i-1}D_{i-1}^*
\]
are elliptic under the first-order elliptic-complex assumptions. Elliptic regularity, compactness of Sobolev inclusions and the spectral theory of elliptic operators supply finite-dimensional smooth harmonic spaces and closed ranges. They then yield the analytic Hodge decomposition and identify harmonic sections with cohomology. These are analytic theorems, not consequences of finite-dimensional linear algebra alone.\cite{v2AS1968I,v2AB1967}

Once those analytic inputs are available, the formal proof of the decomposition follows the same orthogonality identities as in the finite-dimensional case. This division is useful for formalization. A Lean development can first prove the finite-dimensional algebraic theorem, then express the analytic theorem through explicit hypotheses about domains, adjoints, closed range and compactness. It must not suppress these hypotheses merely because the displayed identities look identical.

On noncompact manifolds the situation is more delicate. Continuous spectrum can occur, ranges may fail to be closed, and an operator may not be Fredholm. On manifolds with boundary, a boundary condition must be specified. The closed-manifold index statement below does not automatically extend to either case. Since the Yang--Mills mass-gap problem concerns a theory on four-dimensional Euclidean space and its reconstructed quantum dynamics, this distinction is central rather than peripheral.

\section{The index as an additive invariant}

An operator $D$ between appropriate Hilbert spaces is Fredholm if its kernel and cokernel are finite dimensional and its range is closed. Its index is
\[
\operatorname{ind}D=\dim\ker D-\dim\operatorname{coker}D.
\]
For an elliptic complex, assemble $D+D^*$ from the even-degree to the odd-degree bundles. The analytic Hodge theorem identifies its index with
\[
\chi(E_\bullet,D_\bullet)=\sum_i(-1)^i\dim H^i(E_\bullet,D_\bullet).
\]
The symbol sequence defines a compactly supported topological K-class on $T^*X$, because it is exact away from the zero section.

\begin{theorem}[Atiyah--Singer index theorem; external theorem]
For a complex elliptic differential operator on a closed smooth manifold, the analytic index equals the topological index of its principal-symbol class:
\[
\operatorname{ind}_{\mathrm{an}}(D)
=\operatorname{ind}_{\mathrm{top}}([\sigma(D)]),
\qquad [\sigma(D)]\in K_c^0(T^*X).
\]
The topological index is defined using the Thom construction, an embedding into Euclidean space and Bott periodicity, and is independent of the choices.\cite{v2AS1968I}
\end{theorem}

The theorem turns analytic finite-dimensional defects into a computable topological invariant. It does not determine all eigenvalues, establish a uniform coercive estimate, or reconstruct an interacting quantum field theory. The class $[\sigma(D)]$ contains symbol information; an arithmetic class in $K_0(\Lambda)$ reaches it only after an explicitly constructed geometric map. This is the exact place where the projective-module discussion from Roggenkamp can enter an index-theoretic construction.

A useful specialization is the Dolbeault complex of a holomorphic vector bundle $E$ on a compact complex manifold $X$. Its index is the holomorphic Euler characteristic. The cohomological index formula is
\[
\chi(X,E)=\int_X\operatorname{ch}(E)\operatorname{Td}(T^{1,0}X).
\]
This formula is stated here as a cited external theorem. It should not be confused with a proof that every arithmetic elliptic curve is canonically the same as a chosen complex torus.\cite{v2AS1968III}

\begin{example}[An elliptic curve over the complex numbers]
Let $X$ be a complex elliptic curve. Its holomorphic tangent line bundle is trivial, so its Todd class is one. For a holomorphic line bundle $L$ of degree $d$, the index formula gives
\[
h^0(X,L)-h^1(X,L)=d.
\]
For $L=\mathcal O_X$, the two dimensions are both one and the index is zero. This is an index-theoretic calculation on a genus-one curve. It says nothing by itself about the difficulty of a discrete logarithm in $E(\mathbb F_q)$, which is a finite-group problem with a different input model.
\end{example}

\section{A complete counterexample to a uniform-gap inference}

Consider the circle of length $L>0$, with coordinate $x$ modulo $L$. The functions
\[
f_n(x)=\exp(2\pi\sqrt{-1}nx/L),\qquad n\in\mathbb Z,
\]
are eigenfunctions of $-d^2/dx^2$ with eigenvalues
\[
\lambda_n=4\pi^2n^2/L^2.
\]
The smallest positive eigenvalue is $4\pi^2/L^2$. Both the de Rham complex and its index are independent of $L$ up to the relevant topological equivalence: the Euler characteristic of the circle is zero. Yet its positive spectral gap approaches zero as $L\to\infty$.

To verify that no smaller positive eigenvalue is missing, let $-f^{\prime\prime}=\lambda f$ with periodic values and derivatives. Integration by parts gives $\lambda\|f\|^2=\|f^{\prime}\|^2$, so $\lambda\geq0$. For $\lambda>0$, the elementary ODE solution is a linear combination of $\cos(\sqrt\lambda x)$ and $\sin(\sqrt\lambda x)$. A nonzero periodic solution exists exactly when rotation through $\sqrt\lambda L$ fixes a nonzero vector, equivalently $\sqrt\lambda L\in2\pi\mathbb Z$. This proves the stated least positive eigenvalue. The cohomology has dimension one in degrees zero and one: closed zero-forms are constant, and a one-form is exact precisely when its integral around the circle vanishes. The example gives a rigorous reason that an index theorem cannot alone supply the kind of scale-uniform lower bound needed in a continuum mass-gap argument. The missing ingredient is quantitative control under a limiting process, not another application of Euler-characteristic cancellation.

\section{Connections and curvature: when a differential squares to zero}

Let $P\to X$ be a principal bundle for a compact Lie group $G$, and let $\operatorname{ad}P$ be its adjoint bundle. A connection $A$ defines the covariant exterior derivative
\[
d_A:\Omega^i(X,\operatorname{ad}P)\longrightarrow
       \Omega^{i+1}(X,\operatorname{ad}P).
\]
In a local trivialization one can write $d_A\alpha=d\alpha+[A\wedge\alpha]$. Its curvature is $F_A=dA+\tfrac12[A\wedge A]$.

\begin{proposition}[Curvature identity]
For an adjoint-valued differential form $\alpha$,
\[
d_A^2\alpha=[F_A\wedge\alpha].
\]
\end{proposition}
\begin{proof}
Expand $d_A(d_A\alpha)$ using the graded Leibniz rule for exterior differentiation and the graded Lie bracket. The terms containing one ordinary derivative of $\alpha$ cancel. The remaining terms are $[dA\wedge\alpha]+\tfrac12[[A\wedge A]\wedge\alpha]$, where the graded Jacobi identity supplies the factor $\tfrac12$. This is precisely $[F_A\wedge\alpha]$.
\end{proof}

If $A$ is flat, meaning $F_A=0$, the covariant de Rham sequence is a complex. Its principal symbol is ordinary wedge multiplication by the covector tensored with the identity on $\operatorname{ad}P$. The previously proved Koszul exactness therefore makes it elliptic. Flatness is a sufficient condition for the complex identity; in the adjoint representation a central curvature can also act trivially. It would be inaccurate to say that $d_A^2=0$ always forces the entire curvature to vanish for an arbitrary structure group.

For a general Yang--Mills connection, the equation is $d_A^*F_A=0$. This does not imply $F_A=0$. Thus the full covariant de Rham sequence is not automatically a complex at a general Yang--Mills connection. Writing down $d_A$ in successive degrees and calling it an elliptic complex without checking this composition is a mathematical error.

\section{The anti-self-dual deformation complex}

Assume now that $X$ is an oriented Riemannian four-manifold. On two-forms the Hodge star satisfies $*^2=1$, so
\[
\Omega^2=\Omega^{2,+}\oplus\Omega^{2,-},
\qquad P_+=\tfrac12(1+*).
\]
A connection is anti-self-dual (ASD) when $F_A^+=0$. Set $d_A^+=P_+d_A$ on one-forms. At an ASD connection there is a complex
\[
0\longrightarrow\Omega^0(\operatorname{ad}P)
 \xrightarrow{d_A}\Omega^1(\operatorname{ad}P)
 \xrightarrow{d_A^+}\Omega^{2,+}(\operatorname{ad}P)
 \longrightarrow0.
\]

\begin{proposition}[Complex identity at an ASD connection]
If $F_A^+=0$, then $d_A^+d_A=0$ on adjoint-valued zero-forms.
\end{proposition}
\begin{proof}
The curvature identity gives $d_A^+d_A\phi=P_+[F_A,\phi]$. The Hodge projection acts only on the form component and commutes with the pointwise Lie bracket with the zero-form $\phi$. Hence $P_+[F_A,\phi]=[F_A^+,\phi]=0$.
\end{proof}

\begin{proposition}[Ellipticity of the ASD complex]
The ASD deformation complex is elliptic.
\end{proposition}
\begin{proof}
It suffices to check a fiber and a nonzero covector $\xi$. Rescale and choose an oriented orthonormal basis with $\xi=e^1$. Suppress the adjoint factor, which is unchanged by the symbol maps. The first symbol sends a scalar to its multiple of $e^1$ and is injective. The second sends
\[
\alpha=a_1e^1+a_2e^2+a_3e^3+a_4e^4
\]
to
\[
\frac12\left[
 a_2(e^{12}+e^{34})
 +a_3(e^{13}-e^{24})
 +a_4(e^{14}+e^{23})\right].
\]
The three displayed two-forms are a basis of the self-dual two-forms. Thus the second symbol is surjective and its kernel is the span of $e^1$, exactly the image of the first. Tensoring with the adjoint fiber preserves exactness.
\end{proof}

These are complete proofs of the algebraic identities underlying the deformation complex. The subsequent global study of its moduli space requires elliptic analysis and nonlinear geometry. Atiyah, Hitchin and Singer's work supplies the classical four-dimensional framework and its index calculations.\cite{v2AHS1978}

To see why this complex describes deformations, perturb the connection by a one-form $a$. The curvature expansion is
\[
F_{A+a}^+=F_A^++d_A^+a+\tfrac12[a\wedge a]^+.
\]
At an ASD background the linearized equation is $d_A^+a=0$. An infinitesimal gauge transformation generated by a zero-form $\phi$ changes $a$ by $d_A\phi$, up to the conventional overall sign. Thus infinitesimal deformations modulo infinitesimal gauge are
\[
H_A^1=\ker d_A^+/\operatorname{im}d_A.
\]
The group $H_A^0$ measures infinitesimal stabilizers, and $H_A^2$ is the obstruction space appearing after solving the nonlinear equation locally. The names describe their roles; the statement that a local moduli space has a given smooth dimension requires the stabilizer and obstruction hypotheses below.

\section{Expected dimension and actual local dimension}

Introduce the gauge-fixed operator
\[
\mathcal D_A=d_A^*\oplus d_A^+:
\Omega^1(\operatorname{ad}P)\to
\Omega^0(\operatorname{ad}P)\oplus\Omega^{2,+}(\operatorname{ad}P).
\]
It is elliptic. On a closed manifold its kernel is the harmonic representative space of $H_A^1$. Its cokernel is $H_A^0\oplus H_A^2$. The latter assertion follows from the adjoint formula and orthogonality: the images of $d_A$ and $(d_A^+)^*$ are orthogonal because $d_A^+d_A=0$. Therefore
\[
\operatorname{ind}\mathcal D_A
=\dim H_A^1-\dim H_A^0-\dim H_A^2.
\]

For an $SU(2)$ bundle on a connected $X$ with second Chern number $k$, the standard index calculation is
\[
\operatorname{ind}\mathcal D_A
=8k-3\bigl(1-b_1(X)+b_2^+(X)\bigr).
\]
This formula is cited as an external index-theoretic result. If the connection is irreducible and regular, so that the relevant infinitesimal stabilizer and obstruction spaces vanish, elliptic slice and implicit-function theorems identify the local moduli dimension with this index. Without those hypotheses, it is an expected dimension, not an unconditional dimension of a smooth manifold.\cite{v2AHS1978}

This is a particularly clear instance of the common pattern across the manuscript. A differential identity gives a complex. Its cohomology controls an infinitesimal classification problem. An index evaluates an alternating sum. Nonlinear existence, global classification and quantitative estimates remain additional questions. Replacing the last three questions by the alternating sum would discard information rather than solve them.

\section{Classical action bounds do not establish a quantum mass gap}

Choose an invariant positive inner product on the compact Lie algebra. In four dimensions decompose the curvature as $F_A=F_A^++F_A^-$. Orthogonality gives
\[
\|F_A\|_{L^2}^2=\|F_A^+\|_{L^2}^2+\|F_A^-\|_{L^2}^2.
\]
The quadratic topological expression obtained by pairing $F_A$ with $*F_A$ is
\[
Q(A)=\|F_A^+\|_{L^2}^2-\|F_A^-\|_{L^2}^2.
\]
Hence
\[
\|F_A\|_{L^2}^2\geq|Q(A)|,
\]
with equality when one of the two curvature components vanishes. For the usual trace normalization on $SU(2)$, $|Q(A)|=8\pi^2|k|$. The inequality itself follows only from nonnegativity of the two squared norms; the integrality and topological interpretation of $k$ use Chern--Weil theory.

An ASD connection also solves the classical Yang--Mills equation: the Bianchi identity $d_AF_A=0$ and $*F_A=-F_A$ imply $d_A*F_A=0$, which is equivalent to the Yang--Mills equation. These are classical statements about a connection and its action. A quantum mass gap concerns a positive separation above the vacuum in the spectrum of a rigorously constructed quantum theory. A lower bound on classical action in a nontrivial topological sector does not construct that theory, identify its vacuum Hilbert space, or control its spectrum. The action bound is valuable but proves a different statement.

\section{A lower-dimensional bridge through holomorphic bundles}

Atiyah and Bott studied the Yang--Mills functional on a compact Riemann surface using gauge symmetry and Morse theory, connecting its critical behavior with moduli of holomorphic vector bundles. This gives a strong, established link between complex algebraic geometry, gauge theory and topology. It is a result in two real dimensions with its specified geometric setting; it is not a construction of four-dimensional quantum Yang--Mills theory.\cite{v2AB1983}

On a complex curve, a unitary connection determines a $(0,1)$ operator. Since $(0,2)$-forms vanish in complex dimension one, its integrability condition is automatic, so it defines a holomorphic structure. In higher complex dimensions the curvature component $F_A^{0,2}$ need not vanish, and the same step needs an additional condition. This elementary dimensional observation already explains why a surface result cannot simply be reused in four real dimensions.

For an elliptic curve viewed as a complex Riemann surface, the bundle and Dolbeault constructions fit this lower-dimensional bridge. For an elliptic curve over a finite field, Frobenius, endomorphism rings and finite groups supply arithmetic structures that are not encoded by this analytic surface argument. A comparison project can study both sides, but its claims should specify whether a curve, bundle, module, complex, cohomology group or algorithm is being transported. That discipline is what turns a collection of suggestive themes into a mathematically checkable program.

\section{Interfaces suitable for formal verification}

The preceding chapters identify a finite list of statements with a manageable algebraic core. The chain condition $BA=0$ implies image containment in a kernel. Rank--nullity converts that containment into exactness when dimensions add correctly. The derivative of the chain condition is $BU+VA=0$. A graded commutator squares to zero when the underlying differential does. The Euler identity follows from two exact sequences and cancellation. The wedge--contraction relation proves exactness of a symbol sequence. Each assertion can be represented without postulating the Yang--Mills conjecture or the full Atiyah--Singer theorem.

An honest formalization separates these checked statements from analytic imports. A proof of the finite-dimensional Hodge decomposition does not establish elliptic regularity. A polynomial proof of the ASD symbol calculation does not establish the nonlinear slice theorem. A theorem conditional on a Fredholm hypothesis does not prove that a particular unbounded operator has a closed range. A definition of a topological index is not a formal proof of its equality with the analytic index. The accompanying claim manifest should retain these distinctions so that the word ``verified'' refers to an actual kernel-checked declaration and its visible assumptions.

For future work, the natural formalization layers are: finitely supported complexes and polynomial identities; vector-space dimensions and exactness; finite-dimensional adjoints and positivity; projective-module Euler classes; Sobolev and Fredholm interfaces; and finally geometric symbol classes and index comparisons. This order makes dependencies explicit. It also allows a genuinely useful verified core to be completed while leaving the deep analytic theorems marked as external mathematical inputs.
